\documentclass[10pt,a4paper]{amsart}
\usepackage[margin=19mm]{geometry}
\usepackage{amsmath,amssymb,mathtools}
\usepackage[scr=rsfs,cal=euler]{mathalfa}
\usepackage{xfrac}
\usepackage[unicode,hidelinks,bookmarks=false]{hyperref}
\newcommand{\ie}{i.e.\ }
\newcommand{\cf}{cf.\ }
\newcommand{\resp}{respectively\ }
\newcommand{\Subsection}{\subsection}
\providecommand{\nobreakdash}{\nobreak\hbox{-}}
\setlength{\emergencystretch}{2em}
\multlinegap0pt
\usepackage{amssymb}
\usepackage{mathtools}
\usepackage{xcolor}
\newtheorem{theorem}{Theorem}
\def\k{{\varkappa}}
\def\s{{\sigma}}
\title{Panmagic permutations and $N$-ary groups}
\author{Sergiy Koshkin$^*$ and Jaeho Lee$^\dagger$}
\date{Source: arXiv:2606.22221v1 (CC BY 4.0)}
\begin{document}
\maketitle

\noindent\emph{Source and licence.} Panmagic permutations and $N$-ary groups. Version arXiv:2606.22221v1 (CC BY 4.0), distributed by arXiv under the Creative Commons Attribution 4.0 International licence, http://creativecommons.org/licenses/by/4.0/, as recorded in the arXiv licence metadata for this exact version and shown on the abstract page https://arxiv.org/abs/2606.22221v1. Full licence notice: \texttt{licenses/arxiv-2606.22221-CC-BY-4.0.txt}.\par\smallskip

\noindent\textit{Editorial scope.} Counted unit: the article's theorem criteria (source0.tex lines 230--241). The article's complete original proof of it is delivered with it (source0.tex lines 242--253, one \texttt{proof} environment of 12 lines). No mathematics is written, repaired or re-derived: the statement and the proof are the article's own lines, in the article's own notation, and no retained line is substituted or re-flowed.  No further source-local context is needed: every cross-reference of every retained line resolves inside the two delivered spans.  Nothing was omitted from any retained span, so the package carries no omission record.  No retained line contains a citation, so the wrapper carries no bibliography and no bibliography entry.  No retained line contains a figure environment, an image-inclusion command or drawing code, so no image is delivered and the package declares no asset dependency.  Editorial formatting only, in addition: an AMS \texttt{amsart} preamble that restates the article's own declarations and macros used by the retained lines, namely the environments \texttt{theorem} on separate counters, together with the standard packages those lines need.  The retained environments print in this document as Theorem 1 in order of appearance, whereas the article prints the same items with the numbering of its own full document.  The complete native source of the article as distributed in the arXiv e-print for this exact version is bundled unchanged as \texttt{sources/203332/source0.tex}.
\begin{theorem}[Lionnet]\label{pmagic criteria} Let $\s\in S_n$ be a permutation. Then
\medskip

\noindent\textup{(i)} $\s$ is dihedral if and only if $\sigma(i)-i$ or $\sigma(i)+i$ is constant;
\medskip

\noindent\textup{(ii)} $\s$ is magic if and only if  $0$ has a unique pre-image modulo $n$ under $\sigma(i)-i$ and $1$ has a unique pre-image modulo $n$ under $\sigma(i)+i$;
\medskip

\noindent\textup{(iii)} 
$\s$ is panmagic if and only if $i\mapsto\sigma(i)\pm i$ are both injective modulo $n$, and hence also permutations. 
\end{theorem}

\begin{proof}\phantom{a} 

\noindent(i) The ``only if" part is trivial from the modular formulas for $\phi$ and $\k$. Conversely, if $\s(i)-i=c$ then $\s(i)=i+c$, so $\s=\k^c$. And if $\s(i)+i=c$ then $\s(i)=1-i+c-1$, so $\s=\k^{c-1}\phi$. In both cases, $\s$ is dihedral.
\smallskip

\noindent(ii) Let $\s$ be magic. Since it has a single fixed point there is only one $i$ with $\s(i)=i$, so $0$ has a unique pre-image under $\s(i)-i$. Since it also has a single flipped point, there is only one $j$ with $\s(j)=1-j$ and $1$ has a unique pre-image under $\s(j)+j$. Running the argument in reverse, we conclude that $\s$ has one fixed and one flipped point, and hence is magic. 
\medskip

\noindent(iii) Let $\s$ be panmagic. Then $\k^{-k}\s$ is magic for any $k$, and $\k^{-k}\s(i)-i=\s(i)-i-k=0$ has a unique solution modulo $n$ by (ii). But then $\s(i)-i=k$ has a unique solution for any $k$ and $\s(i)-i$ is injective. Similarly, $\k^{k}\phi\s$ is magic for any $k$, $\k^{k}\phi\s(i)-i=1-(\s(i)+i)+k=1$ has a unique solution, and then so does $\s(i)+i=k$. Thus, $\s(i)+i$ is also injective. 

Conversely, let $\sigma(i)\pm i$ be injective modulo $n$. By (ii), it is enough to show that so are $\delta\s(i)\pm i$ for all $\delta\in D_n$, and even just for $\delta=\k,\phi$ since they generate $D_n$. But $\k\s(i)\pm i=\sigma(i)\pm i+1$ and $\phi\s(i)\pm i=1-(\sigma(i)\mp i)$, so this follows trivially.
\end{proof}

\end{document}
